% !TEX program = pdflatex
%
% Bracket-Structured Probability Estimation via Quantum Walks
% and Live Prediction-Market Validation
%
% arXiv preprint, draft v0 — June 2026
% Categories : cs.ET (Emerging Technologies) ; cross-list : quant-ph, q-fin.CP

\documentclass[11pt]{article}
\usepackage[utf8]{inputenc}
\usepackage[T1]{fontenc}
\usepackage{lmodern}            % cleaner body font (vs default CM bitmaps)
\usepackage{microtype}          % subtle kerning, protrusion, expansion
\usepackage[a4paper,
            top=1in, bottom=1in,
            left=1.15in, right=1.15in]{geometry}
\usepackage{amsmath, amssymb, amsthm}
\usepackage{graphicx}
\usepackage{hyperref}
\hypersetup{colorlinks=true, urlcolor=blue, citecolor=blue,
            linkcolor=black}
\usepackage{booktabs}
\usepackage{authblk}
\usepackage{enumitem}
\usepackage{xcolor}
\usepackage{placeins}           % \FloatBarrier keeps floats in their section

% Slightly tighten section spacing — default \section spacing is generous and
% pushes sections to the top of a page with too much air. We keep \subsection
% spacing untouched (the contrast is what flags hierarchy at a glance).
\usepackage{titlesec}
\titlespacing*{\section}{0pt}{2.0ex plus 0.3ex minus 0.2ex}{1.2ex plus 0.2ex}
\titlespacing*{\subsection}{0pt}{1.6ex plus 0.3ex minus 0.2ex}{0.8ex plus 0.2ex}

% Avoid orphan / widow lines at page breaks — 10000 is "never".
\widowpenalty=10000
\clubpenalty=10000

% Slightly looser caption / float spacing for readability.
\setlength{\textfloatsep}{14pt plus 4pt minus 2pt}
\setlength{\intextsep}{12pt plus 3pt minus 2pt}

\title{Bracket-Structured Probability Estimation via Quantum Walks\\
and Live Prediction-Market Validation}

\author{Jean-Thomas Dauchel\thanks{\texttt{contact@pulsquant.com}}}

\affil{PulsQuant --- Independent Research}

\date{June 2026}

\begin{document}
\maketitle

\begin{abstract}
This paper applies quantum walks to sports tournament brackets. To
our knowledge it is the first formal treatment of the idea : the 2018
\emph{Quantum Frontiers} note \cite{quantumfrontiers_wc_2018} sketched
a qudit encoding for the FIFA World Cup but never went through peer
review, and the recent Nature/QNN football work
\cite{nature_qnn_football_2025} targets a different task (match-level
classification by QML).

We map an $n$-team single-elimination bracket to a circuit of $n{-}1$
qubits with controlled-$R_Y$ gates. Round-0 matches use the exact
Elo logistic ; deeper rounds use a \emph{seeded} approximation whose
bias we quantify against a closed-form enumeration that follows the
same matchup rule, so any $Q$--$C$ divergence is attributable to the
approximation rather than to a quantum implementation artifact.

The encoding is back-tested on five real 2025--2026 brackets loaded
from a production prediction-market database (not balanced
assumptions). On the 2026 NBA Playoffs --- where the San Antonio
Spurs reached the NBA Finals from the 8-seed in the West, eliminating
the 1-seed OKC in the Western Conference Finals --- the seeded walk
assigns $P_Q(\text{SAS reaches Finals}) = 5.87\%$ against $P_C =
3.43\%$ ($\mathbf{1.71\times}$ upweight) and a $\mathbf{3.05\times}$
upweight on $P(\text{SAS champion})$. On the Roland Garros WTA bracket
from the Round of 16 onward, the eventual champion Mirra Andreeva is
upweighted $\mathbf{1.68\times}$ on her champion marginal and the
underdog finalist Maja Chwalinska (Tennis Elo 1710, WTA \#114
pre-tournament) is upweighted $\mathbf{2.78\times}$ on her
runner-up marginal. Underdog semifinalists on the ATP side (Jakub
Mensik, Matteo Arnaldi) come out at $\mathbf{4.17\times}$ and
$\mathbf{4.44\times}$. Two clean negative controls confirm the
method is not just an upset-amplifier : the chalk-favourite PSG
winning the UCL gives $Q/C \approx 1.01$, and Vegas sweeping Colorado
in the NHL Stanley Cup Conference Finals --- statistically a coin
flip despite the sport-narrative upset framing --- gives $Q/C =
0.63\times$ on Vegas reaching the Final.

As operational extensions of the encoding, we implement both Grover
amplitude amplification and Quantum Amplitude Estimation
\cite{brassard2002quantum} on the SAS-champion query of the NBA
bracket. The amplification curve peaks at $97.89\%$ on $k = 4$
iterations (textbook optimum), and the full QAE with $m = 6$ ancilla
qubits (21 qubits total on a single-laptop Aer simulator) returns
$\hat p_{\text{mode}} = 3.81\%$ on the SAS-champion amplitude versus
the reference $p_0 = 3.46\%$ from unamplified sampling. The QAE
mode and the weighted-mean estimator agree with the reference within
the $\pi / 2^m$ resolution band at every $m \in \{4, 5, 6\}$,
realising the quadratic Brassard speedup on a real-bracket query.

A prospective prediction for the FIFA World Cup 2026 (kickoff
June 11, 2026) is issued in the same form : mid-tier nations (NED,
BEL, AUT, COL) upweighted by 2--3 percentage points, top European
favourites (FRA, ESP, BRA) commensurately reduced. ARG is the notable
exception among top favourites --- upweighted $+6.14$\,pp because its
balanced-draw path runs through a softer Conmebol-heavy quarter.

At simulator scales (15 qubits) we make no claim of quantum advantage.
The contribution is conceptual : the encoding captures bracket
structural correlations that classical Monte Carlo only reproduces
via explicit path-by-path conditioning, and the construction scales
tractably toward the 50--100 logical qubits expected from early
fault-tolerant hardware.
\end{abstract}

\section{Introduction}
\label{sec:intro}

Sports tournaments structured as single-elimination brackets exhibit
strong \emph{structural correlations} between match outcomes : if a
high-seed team loses early, downstream matches are reshuffled in
ways that affect the marginal champion probability of every other
participant. Pricing these correlations correctly is a central
problem for prediction markets and bookmakers, who must aggregate
expectation over an exponential number of bracket paths to emit
champion futures odds. Classical Monte Carlo (MC) is the standard
approach \cite{glasserman2004monte}, with the textbook estimator
running $N$ independent simulations of the full bracket and counting
champion outcomes \cite{simulate_bracket_pulsquant}. Each simulation is
$O(n)$ in the number of matches $n$, and Monte-Carlo sample noise on
the marginal of a low-probability champion scales as $O(1/\sqrt{N})$.
On the match-level modelling side, the closest classical references
are the Dixon-Coles family of Poisson-based football models
\cite{dixon1997modelling}, and their Bayesian extensions
\cite{constantinou2017bayesian} : these set the implied match
probabilities, while the structural correlation between matches is
typically left to the MC layer.

The 2025--2026 season delivered an unusually rich set of bracket
upsets to test such an estimator against. In the NBA Playoffs the
San Antonio Spurs reached the NBA Finals from the 8-seed in the
West, eliminating the 1-seed OKC Thunder in the Western Conference
Finals. Roland Garros 2026 lost three of the top-four seeds before
the quarterfinals (Sinner on the men's side, Świątek and Sabalenka
on the women's). In the NHL Stanley Cup the Vegas Golden Knights
eliminated the regular-season-best Colorado. None of these results
is statistically extreme on its own match marginal --- a well-
calibrated Elo gives each of them double-digit probability. What
makes them difficult to price is the joint weight they carry in the
bracket tree once a deep run is conditioned on : the correlation
between successive rounds amplifies the marginal of the surviving
underdog far above what an independent-rounds approximation would
suggest.

This paper proposes a quantum walk that encodes those correlations
through entanglement between conditionally-dependent match qubits,
and studies its behaviour in three settings : an 8-team synthetic
bracket as a sanity check ; five real 2025--2026 tournaments
back-tested against ground truth (NBA Playoffs, NHL Stanley Cup
Conference Finals, UEFA Champions League, and the men's and women's
Roland Garros draws from the Round of 16 onward) ; and the 2026 FIFA
World Cup R16-onward (15 matches), with the prediction issued before
the June~11 kickoff. All brackets are loaded from the production
PulsQuant prediction-market database rather than from balanced
toy assumptions.

The construction picks up where two recent threads left off. The
2018 \emph{Quantum Frontiers} note \cite{quantumfrontiers_wc_2018}
sketched a qudit encoding for the World Cup bracket but neither
published the formalism nor implemented the Grover path-likelihood
queries it flagged as future work. The Nature/QNN football paper
\cite{nature_qnn_football_2025} attacks a different target : match-
level outcome classification via quantum machine learning, leaving
the bracket-structural problem open. On the classical side,
\cite{brandes2025stop} drops Monte Carlo altogether by exploiting
independence between bracket sub-paths ; we go the opposite way and
exploit the structural \emph{dependence} between adjacent matches.

\section{Background}
\label{sec:bg}

\subsection{Quantum walks}

A discrete-time quantum walk on a graph $G$ replaces classical
transition matrices with unitary operators acting on the tensor
product of a position Hilbert space and a coin Hilbert space
\cite{aharonov1993quantum, kempe2003quantum, portugal2018quantum}.
For a single-elimination bracket viewed as a balanced binary tree
of depth $d$ ($n = 2^d$ teams, $2^d - 1$ internal matches), the
canonical walk traverses the tree level-by-level toward the root,
with the coin determining the choice of winning child at each
match. The exponential algorithmic speedups of Childs, Cleve,
Deotto, and Farhi \cite{childs2003exponential} apply to a different
problem (hidden-graph traversal) ; we keep the structural encoding
but make no asymptotic claim of that flavour here.

\subsection{Quantum amplitude estimation}

For a state $|\psi\rangle$ that places weight $p$ on a target
subspace, classical sampling requires $O(1/p\epsilon^2)$ shots to
estimate $p$ to additive error $\epsilon$. The amplitude estimation
algorithm of Brassard, Hoyer, Mosca, and Tapp
\cite{brassard2002quantum} achieves the same precision in
$O(1/p\epsilon)$ queries, a quadratic speedup which Stamatopoulos
\textit{et al.} \cite{stamatopoulos2020option} exploited for the
end-to-end pricing of European call options on near-term IBM hardware.
Subsequent refinements of QAE that remove the phase-estimation
register \cite{suzuki2020amplitude} cut the resource overhead by
several factors, making the technique tractable on NISQ devices
\cite{preskill2018nisq, bharti2022nisq}. A parallel line of work
extends the same machinery to Monte Carlo speedups more broadly
\cite{montanaro2015quantum}, with credit-risk pricing as one of the
early industrial demonstrations \cite{egger2020option}.

\subsection{Quantum-enhanced finance, 2024--2026}

The most directly comparable recent industrial benchmark is the
Stamatopoulos et al. \cite{stamatopoulos2020option} pipeline for
end-to-end pricing of European options via QAE on near-term IBM
hardware, subsequently scaled by the IBM Quantum team to its
127-qubit superconducting devices in the 2023--2024 NISQ era
\cite{bharti2022nisq}. The Goldman--IBM
resource estimate \cite{goldman_ibm_2021} pegged true quantum
advantage in derivative pricing at $\sim 7500$ logical qubits and
T-depth $\sim 4.6 \times 10^7$, a regime still beyond any current
hardware but well below the historical estimates of the late 2010s.

\section{Method}
\label{sec:method}

\subsection{Bracket encoding}

Let $T = \{t_1, \dots, t_n\}$ denote the participating teams with
Elo ratings $\rho_i$, and let $\mathcal{M} = \{m_0, \dots, m_{n-2}\}$
denote the matches in topological order (round-0 first). To each
match $m \in \mathcal{M}$ we associate a qubit $q_m$. The convention
$q_m = 0$ encodes ``the team in the upper bracket position of $m$
won'' ; $q_m = 1$ encodes ``the lower team won''.

\paragraph{Round 0.} The participants are known up front. For each
first-round match $m$ pairing teams $a$ and $b$, we compute
$p_{m} = 1/(1 + 10^{(\rho_b - \rho_a)/400})$, the classical logistic
Elo win probability, and apply
\[
R_Y\!\left(\theta_m\right) |0\rangle, \qquad
\theta_m = 2 \arcsin\!\sqrt{1 - p_m},
\]
which leaves $q_m$ in the state $\sqrt{p_m}\,|0\rangle +
\sqrt{1-p_m}\,|1\rangle$, so a computational-basis measurement
returns $|0\rangle$ with the correct marginal.

\paragraph{Deeper rounds (seeded approximation).} For a match $m$ at
round $r > 0$, the actual participants depend on the outcomes of the
two parent matches $\pi^t(m), \pi^b(m)$. A faithful encoding would
condition on all $2^{r-1}$ possible parent configurations using
$r$-qubit-controlled rotations ; this scales as $O(2^d)$ in the
bracket depth. We instead apply a \emph{seeded} approximation : for
each of the 4 combinations of parent outcomes, we identify the team
that would emerge under a canonical seed-following path (always
top-side if the parent bit is 0, always bot-side if 1), compute the
matchup probability under those seed teams, and apply the
corresponding multi-controlled $R_Y$ rotation. The seeded
approximation is exact for first-round matches and at the
$d$-axis-aligned paths through the bracket ; it is biased toward
those canonical seed matchups for deeper rounds.

\subsection{Closed-form check of the seeded approximation}

Alongside $Q$ (Qiskit Aer sampling) and the unbiased classical
Monte Carlo $C$, we compute a third estimator $S$ that enumerates
all $2^{n-1}$ bracket outcomes in closed form under the \emph{same}
seeded matchup rule used by the circuit. The triple
$\{Q, C, S\}$ then lets us split the $Q$--$C$ divergence into two
parts : whatever sits between $Q$ and $S$ is sampling noise or
implementation error, while whatever sits between $S$ and $C$ is
the seeded approximation's bias. Throughout the experiments we
find $|Q - S| < 0.1$\,pp, so the gap to $C$ is attributable to the
approximation alone.

\section{Experiments}
\label{sec:experiments}

\subsection{Implementation}

All experiments run on Qiskit 2.4.1 + Qiskit Aer 0.13 in
single-laptop CPU mode. We use 50\,000 shots per quantum circuit
and 50\,000 sims per classical MC, with a fixed seed (42) for
reproducibility. Source code and raw data are available at
\url{https://github.com/pulsquant/quantum-bracket}.

\subsection{Sanity check : 8-team bracket}

A 7-match, 8-team bracket with linearly spaced Elo ratings
($\rho_A = 1850, \rho_B = 1800, \dots, \rho_H = 1500$) provides
a low-depth setting where the seeded approximation's bias should
be small. Figure~\ref{fig:fig1} confirms that $Q$ and $C$ agree
to within max $|\Delta| = 1.48$\,pp (limited by 50k-shot sampling
noise) across all 8 candidate champions.

\begin{figure}[!htbp]
\centering
\includegraphics[width=\linewidth]{figures/fig1_smoke_test.pdf}
\caption{Sanity check : 8-team synthetic bracket. Quantum walk (blue)
and classical Monte Carlo (orange) champion distributions agree
within 1.5\,pp across all teams.}
\label{fig:fig1}
\end{figure}

\subsection{Real-bracket back-tests : NBA, NHL, UCL, RG ATP, RG WTA}
\label{sec:multi}

We load the actual 2025--2026 brackets from the PulsQuant
\texttt{tournament\_brackets} production table and construct the
\texttt{Bracket} object directly from the recorded match pairings.
For the NBA Playoffs and UEFA Champions League knockout stage we
load the full bracket. For Roland Garros (128-entry draws) we slice
from the Round of 16 onward to keep the qubit count at 15. Strength
ratings are taken from a manually-curated table mirroring
\texttt{eloratings.net} (NBA, soccer) and Tennis-Elo (ATP, WTA),
listed in Appendix.

The realised outcomes provide ground-truth back-test targets. Table
\ref{tab:multi_tournament} summarises the quantum walk and classical
Monte Carlo marginals on the realised champion, runner-up and
semifinalists across the five tournaments.

\begin{table}[!htbp]
\centering
\caption{Multi-tournament back-test : $P(\text{outcome})$ under the
seeded quantum walk $Q$ versus unbiased classical MC $C$, for
realised champions, runners-up, semifinalists and finalists across
five tournaments. Bold $Q/C$ ratios exceed $1.5\times$ or fall below
$0.6\times$. All tournaments use the actual production bracket data ;
tennis tournaments are loaded from the Round of 16 onward.}
\label{tab:multi_tournament}
\small
\begin{tabular}{llrrr}
\toprule
Tournament & Realised role / participant & $P_Q$ & $P_C$ & $P_Q/P_C$ \\
\midrule
NBA Playoffs 2026 & SAS reaches Finals (8-seed, beat OKC WCF) & 5.87\% & 3.43\% & $\mathbf{1.71\times}$ \\
                  & SAS champion (Finals 2--0 down)            & 3.46\% & 1.13\% & $\mathbf{3.05\times}$ \\
                  & NYK reaches Finals (East favourite)        & 13.22\% & 18.75\% & 0.71$\times$ \\
                  & NYK champion (Finals 2--0, ongoing)        & 4.83\% & 9.59\% & $\mathbf{0.50\times}$ \\
                  & OKC top-seed (eliminated WCF)              & 18.52\% & 17.73\% & 1.04$\times$ \\
\midrule
NHL Stanley Cup 2026 & VGK reaches Final (swept COL WCF 4--0) & 12.65\% & 20.00\% & $\mathbf{0.63\times}$ \\
                     & CAR reaches Final (East CF winner) & 26.93\% & 23.60\% & 1.14$\times$ \\
                     & COL top-Elo West (eliminated WCF) & 18.32\% & 19.13\% & 0.96$\times$ \\
                     & OTT R1 loser (low-Elo) & 5.24\% & 2.67\% & $\mathbf{1.96\times}$ \\
\midrule
UCL 2025--26      & PSG champion (top favourite) & 10.57\% & 10.43\% & 1.01$\times$ \\
                  & ARS runner-up                 & 19.00\% & 13.77\% & 1.38$\times$ \\
                  & B04 SF                        & 10.41\% & 4.50\% & $\mathbf{2.32\times}$ \\
                  & RMA (eliminated QF)           & 10.10\% & 14.64\% & 0.69$\times$ \\
\midrule
RG ATP 2026 (R16+) & Zverev finalist (top favourite) & 4.05\% & 38.15\% & $\mathbf{0.11\times}$ \\
                   & Cobolli finalist (underdog)     & 6.27\% & 5.74\% & 1.41$\times$ \\
                   & Mensik SF (underdog)            & 5.59\% & 1.34\% & $\mathbf{4.17\times}$ \\
                   & Arnaldi SF (underdog)           & 3.07\% & 0.69\% & $\mathbf{4.44\times}$ \\
\midrule
RG WTA 2026 (R16+) & \textbf{Andreeva champion (18 y/o)} & \textbf{11.01\%} & \textbf{6.55\%} & $\mathbf{1.68\times}$ \\
                   & \textbf{Chwalinska runner-up (Elo 1710)} & \textbf{2.15\%} & \textbf{0.77\%} & $\mathbf{2.78\times}$ \\
                   & Shnaider SF (eliminated Sabalenka QF) & 3.65\% & 2.49\% & 1.46$\times$ \\
                   & Sabalenka (top favourite, eliminated QF) & 54.79\% & 41.08\% & 1.33$\times$ \\
                   & Swiatek (top favourite, eliminated R16) & 3.74\% & 37.01\% & $\mathbf{0.10\times}$ \\
\bottomrule
\end{tabular}
\end{table}

A consistent pattern emerges. Realised underdog deep runs are
upweighted : the Spurs reaching the NBA Finals as an 8-seed after
eliminating 1-seed OKC ($1.71\times$ on the reach-Finals marginal,
$3.05\times$ on the champion marginal), Andreeva's RG WTA title
($1.68\times$), Chwalinska's runner-up run with a Tennis Elo of
1710 and a WTA ranking of 114 pre-tournament ($2.78\times$), and
the underdog ATP semifinalists Mensik and Arnaldi ($4.17\times$
and $4.44\times$ respectively). Top favourites that lost early in
their bracket are down-weighted in turn, with Świątek the cleanest
case at $0.10\times$ on the WTA side (eliminated R16 by Andreeva).
On the same side, top favourites whose Elo edge is robust against
their actual draw remain at $Q/C \approx 1$ : OKC at $1.04\times$
on the NBA champion marginal, despite being eliminated by SAS in
the WCF, simply because OKC's Elo (1745) was the highest in the
field and the seeded approximation reads its path through the West
as competitive. When the chalk holds outright, the walk returns
essentially the classical value : PSG winning the UCL gives
$Q/C \approx 1.01$, evidence that the encoding does not manufacture
upset signal in its absence.

The price for that asymmetry shows up on top seeds that survive
deep. The New York Knicks reached the NBA Finals from the East,
yet $Q$ assigns their reach-Finals marginal $P_Q = 13.22\%$ versus
$P_C = 18.75\%$ ($0.71\times$), and their champion marginal
$P_Q = 4.83\%$ versus $P_C = 9.59\%$ ($0.50\times$). The seeded
approximation biases against top-seed bracket positions whenever
their Elo lead is moderate enough to be eroded by repeated
conditional matchups ; the upweight on underdog paths and the
down-weight on chalk paths are two sides of the same operation.
The NHL Vegas--Colorado West Final is the cleanest control we
have : narratively an upset, statistically a coin flip (Elo 1645
vs 1670), and the walk does not produce upset signal there either
($Q/C = 0.63\times$ on Vegas reaching the Stanley Cup Final).

The closed-form seeded estimator $S$ produces values within
sampling noise of $Q$ on all five tournaments (max $|Q-S| <
0.09$\,pp), confirming the $Q$--$C$ divergence is the seeded
approximation's bias, not a quantum implementation artifact.

Figure~\ref{fig:fig2} visualises the NBA Playoffs 2026 champion
marginal distribution side-by-side for $Q$ and $C$, illustrating the
characteristic redistribution of probability mass from top seeds
toward bracket-bottom paths.

\begin{figure}[!htbp]
\centering
\includegraphics[width=\linewidth]{figures/fig2_nba_2026.pdf}
\caption{2026 NBA Playoffs champion distribution (15 qubits). Quantum
walk redistributes probability mass selectively : mid-East seeds NYK
and DEN are down-weighted (their classical paths are heavy on
moderate-Elo opponents that the seeded approximation reads as
upset-vulnerable), while the 1-seeds OKC and BOS remain close to their
classical values (their Elo lead is robust enough against the actual
draw that no down-weighting is applied). SAS, the 8-seed West, made
the bracket-relevant deep run of the 2026 playoffs by eliminating OKC
in the Western Conference Finals to reach the NBA Finals --- exactly
the kind of structurally correlated upset path the seeded quantum
walk is designed to upweight ex ante. The NBA Finals series was
ongoing at the time of writing (NYK leading 2--0 in the best-of-7).}
\label{fig:fig2}
\end{figure}

\subsection{Grover amplification on a bracket query}
\label{sec:grover}

A natural follow-up to the encoding is to amplify the marginal of a
specific tournament outcome using Grover's algorithm
\cite{grover1996fast}, rather than estimating it by direct sampling.
The 2026 NBA bracket gives a clean test bed : SAS reaching the NBA
Finals as an 8-seed corresponds to a 4-qubit constraint on the
15-qubit register (the bits indexing the OKC-SAS R1 match, the West
CSF, the West CF, and the NBA Finals being on the West side). We
encode that constraint as a phase oracle $O$ and pair it with the
diffusion operator $D = A (2|0\rangle\langle 0| - I) A^\dagger$, where
$A$ is the bracket prep unitary of Section~\ref{sec:method}.

Figure~\ref{fig:fig4} shows the empirical amplification curve : the
SAS-champion measurement frequency, sampled at $50\,000$ shots per
point, against the number $k$ of Grover iterations applied to $A|0\rangle$.
The baseline value $p_0 = 3.46\%$ matches the conditional encoding's
prediction at $k=0$. Successive iterations follow the textbook
envelope $\sin^2((2k+1)\theta)$ with $\sin\theta = \sqrt{p_0}$, peaking
at $97.89\%$ on $k = 4$ (the theoretical optimum being
$k^\star = \pi / (4 \arcsin\sqrt{p_0}) \approx 4.20$), decaying to
$\sim 1\%$ at $k = 8$, then peaking again at $99.72\%$ on $k = 12$.

The shot-count implication is immediate. Targeting a $5\%$ relative
precision on the SAS-champion frequency, direct sampling requires
$N_{\text{unamp}} = (1 - p_0)/(p_0 \rho^2) \approx 11\,200$ shots,
while sampling from the amplified state with $k = 4$ requires
$N_{\text{amp}} \approx 9$ shots (and $\approx 1$ shot at $k = 12$).
The roughly $10^4 \times$ empirical shot-count reduction is consistent
with the quadratic Grover speedup $O(1/\sqrt{p_0})$ in the
amplitude-amplification sense \cite{brassard2002quantum}. The cost is
circuit depth : each Grover iteration adds one application of
$A$, $O$, and $A^\dagger$, scaling the depth by roughly $2k + 1$.
On 15-qubit Aer this remains tractable on a single laptop CPU, with
the full $k = 0\dots15$ sweep completing in a few minutes.

Amplification rotates the bracket state until the SAS-champion
subspace dominates the measurement, but it does not directly recover
the value of $p_0$ itself : it just rotates it close to 1. Section
\ref{sec:qae} closes that gap with the canonical Brassard-Hoyer-Mosca-
Tapp amplitude estimation procedure on the same Grover operator.

\begin{figure}[!htbp]
\centering
\includegraphics[width=0.85\linewidth]{figures/fig4_grover_speedup.pdf}
\caption{Grover amplification on the SAS-champion query of the 2026
NBA Playoffs bracket. Blue : Aer measurement at $50\,000$ shots per
point. Gray dashed : theoretical envelope $\sin^2((2k+1)\theta)$ with
$\sin\theta = \sqrt{p_0}$ and $p_0 = 3.46\%$ (orange dotted).
Two peaks at $k = 4$ ($97.89\%$) and $k = 12$ ($99.72\%$), consistent
with the textbook Grover periodicity.}
\label{fig:fig4}
\end{figure}

\subsection{Quantum amplitude estimation on the same query}
\label{sec:qae}

The canonical extension of the previous section's Grover amplification
is Quantum Amplitude Estimation (QAE) \cite{brassard2002quantum}, which
returns an estimate of $p_0$ itself rather than amplifying the target
subspace to peak. The construction adds $m$ ancilla qubits to the
$n_{\text{data}} = 15$-qubit bracket register and proceeds in three
steps : (i) the bracket prep $A$ is applied to the data register,
(ii) for $j = 0, \dots, m - 1$ the ancilla $j$ controls $Q^{2^j}$
where $Q$ is the same Grover operator as in
Section~\ref{sec:grover}, with the canonical sign convention so its
eigenvalues on the good/bad subspace are exactly $e^{\pm 2i\theta}$
with $\sin\theta = \sqrt{p_0}$ ; (iii) an inverse QFT on the ancillas
followed by a measurement returns an integer $y$ which encodes the
amplitude estimate $\hat p = \sin^2(\pi y / 2^m)$ on a $2^m$-grid.

We sweep $m \in \{4, 5, 6\}$, totalling 19, 20 and 21 qubits
respectively, which remains tractable on a single-laptop Aer
simulator (each $m$ runs in a few minutes). 10\,000 shots per
configuration. Table~\ref{tab:qae} summarises the convergence to
the reference value $p_0 = 3.46\%$ established by the unamplified
sampling baseline. The mode of the ancilla measurement saturates
to $\hat p = 3.81\%$ for all three values of $m$ : the discrete
$2^m$-grid happens to place its closest point at $y = 15 \cdot
2^{m-4}$ with $\sin^2(\pi \cdot 15/16) \approx 3.81\%$, which is
within the theoretical $\pi / 2^m$ resolution band of the true
value at every $m$. The weighted-mean estimator $\hat p_{\text{mean}}
= \sum_y P(y) \sin^2(\pi y / 2^m)$ refines further with $m$ : $3.84\%$
at $m = 4$, $3.90\%$ at $m = 5$, $3.95\%$ at $m = 6$, all within
$0.5$\,pp of the reference.

\begin{table}[!htbp]
\centering
\caption{QAE convergence on the SAS-champion query. Reference
$p_0 = 3.46\%$ from $50\,000$-shot unamplified sampling
(Figure~\ref{fig:fig4}, $k = 0$). $\hat p_{\text{mode}}$ : mode of
the ancilla measurement histogram, decoded via $\sin^2(\pi y / 2^m)$.
$\hat p_{\text{mean}}$ : empirical-frequency-weighted estimator. The
quadratic Brassard speedup manifests as the
$O(1/\sqrt{p_0}) \approx 5$ controlled-$Q$ iterations sufficient at
$m = 6$, versus the $O(1/p_0) \approx 30$ unamplified samples needed
for the same expected count.}
\label{tab:qae}
\small
\begin{tabular}{cccccc}
\toprule
$m$ & Total qubits & Resolution $\pi / 2^m$ & Peak count $/ 10^4$ &
$\hat p_{\text{mode}}$ & $\hat p_{\text{mean}}$ \\
\midrule
4 & 19 & $0.196$ & $5\,037$ & $3.81\%$ & $3.84\%$ \\
5 & 20 & $0.098$ & $4\,943$ & $3.81\%$ & $3.90\%$ \\
6 & 21 & $0.049$ & $4\,792$ & $3.81\%$ & $3.95\%$ \\
\bottomrule
\end{tabular}
\end{table}

Figure~\ref{fig:fig5} shows the empirical distribution of $\hat p$
at $m = 6$. The mode at $y = 60$ (folded to $y = 4$ by the
$y \leftrightarrow 2^m - y$ symmetry) carries $47.9\%$ of the shots
and corresponds to $\hat p = 3.81\%$ ; the next-largest peaks sit at
$y = 28$ and $y = 36$, which decode to amplitudes around $90\%$
but only collect $<1\%$ of shots each, a sampling tail consistent
with the QPE precision tail \cite{brassard2002quantum,
suzuki2020amplitude}.

\begin{figure}[!htbp]
\centering
\includegraphics[width=0.85\linewidth]{figures/fig5_qae_estimation.pdf}
\caption{QAE empirical estimator distribution at $m = 6$ ancillas
($21$ total qubits, $10\,000$ shots). Blue : histogram of decoded
amplitudes $\hat p = \sin^2(\pi y / 2^m)$. Orange dashed : reference
$p_0 = 3.46\%$ from unamplified sampling. Red dotted : QAE mode
estimate $\hat p_{\text{mode}} = 3.81\%$.}
\label{fig:fig5}
\end{figure}

The shot-count comparison with the amplification scheme of
Section~\ref{sec:grover} clarifies the practical difference between
the two methods. Amplification rotates $p_0$ to a peak near $1$,
making it cheap to verify whether the query subspace is hit, but
does not return $p_0$ itself : extracting a numerical estimate from
the amplified state still requires a separate measurement campaign
to undo the amplification angle. QAE returns $\hat p$ directly from
a single shot of the ancilla register, at the cost of running
$2^m - 1 \approx 63$ Grover iterations per shot at $m = 6$ versus
the single application of $A$ per unamplified sample. The
asymptotic shot-count saving is the well-known quadratic Brassard
speedup $O(1/p_0)$ vs $O(1/p_0^2)$ for additive precision ; at
$p_0 = 3.5\%$ that is a $\sim 30 \times$ reduction in shots, which
the empirical convergence in Table~\ref{tab:qae} corroborates.

\subsection{Prospective prediction : FIFA World Cup 2026}

The 2026 FIFA World Cup kicks off on June~11 across 16 host cities
in the USA, Canada, and Mexico. We issue a prospective prediction
on the R16-onward stage (15 matches, 15 qubits) by taking the top-16
teams ranked by $P(\text{reach R16})$ from the PulsQuant classical
MC snapshot of 2026-06-06T17:04 and slotting them into a standard
balanced 1-vs-16 bracket. Table~\ref{tab:wc} shows the top
entries of the quantum-walk distribution alongside the corresponding
classical values. Mass moves from the top favourites (FRA $-9.5$\,pp,
ESP $-4.8$\,pp) toward the mid-tier (NED $+3.0$\,pp, BEL $+2.9$\,pp,
AUT $+2.2$\,pp, COL $+2.0$\,pp) --- the same redistribution that
the NBA back-test exhibited, transposed to a draw whose ground
truth is still six weeks away.

\begin{table}[!htbp]
\centering
\caption{FIFA World Cup 2026, R16-onward prospective champion
probability (top 12 entries, sorted by $P_Q$). Actual knockout pairings
will be determined by the group stage and may differ from our balanced
1-vs-16 assumption. The mid-tier upweight pattern (NED, BEL, COL, AUT
all $+2$ to $+3$\,pp) mirrors the NBA back-test of
Section~\ref{sec:multi}.}
\label{tab:wc}
\begin{tabular}{lrrr}
\toprule
Team & $P_Q$ (quantum walk) & $P_C$ (classical MC) & $P_Q - P_C$ \\
\midrule
ARG & 31.06\% & 24.92\% & $+6.14$\,pp \\
BRA & 9.26\% & 13.14\% & $-3.88$\,pp \\
GER & 7.54\% & 6.04\% & $+1.49$\,pp \\
POR & 7.19\% & 8.36\% & $-1.17$\,pp \\
NED & 6.70\% & 3.68\% & $+3.02$\,pp \\
BEL & 6.11\% & 3.18\% & $+2.93$\,pp \\
ESP & 5.36\% & 10.14\% & $-4.78$\,pp \\
FRA & 5.17\% & 14.63\% & $-9.45$\,pp \\
COL & 4.33\% & 2.30\% & $+2.03$\,pp \\
ENG & 4.32\% & 5.97\% & $-1.65$\,pp \\
SUI & 2.80\% & 1.49\% & $+1.31$\,pp \\
AUT & 2.61\% & 0.39\% & $+2.22$\,pp \\
\bottomrule
\end{tabular}
\end{table}

\section{Discussion}
\label{sec:discussion}

\subsection{Scope of the claim}

A word on what is and isn't being claimed. At 15-qubit simulator
scales we are not racing the classical MC : the
quantum sampling runs in $0.16$~s and the classical MC in $0.06$~s
(Qiskit Aer on a single CPU thread, 50\,000 shots and sims). The
contribution is the encoding itself, not its present runtime. The
seeded approximation is in particular a controlled and reversible
source of bias --- the closed-form check in
Section~\ref{sec:method}.2 attributes the entire $Q$--$C$ gap to it,
which means upgrading the encoding (next paragraph) recovers
unbiased MC behaviour up to the qubit budget.

\subsection{What the encoding captures}

What the construction does buy is structural correlation between
adjacent matches for free, by entanglement, without explicit path
conditioning. The empirical effect on the five tournaments we tested
is consistent : every realised underdog deep run we have a
back-test for is upweighted, by factors ranging from $1.71\times$
(SAS reaching the NBA Finals) to $4.44\times$ (Arnaldi as an ATP
underdog semifinalist), and the WTA champion Andreeva sits in the
middle of that range at $1.68\times$. The down-weighting on top
seeds that survive deep --- the New York Knicks reach-Finals
marginal at $0.71\times$ and their champion marginal at
$0.50\times$ being the cleanest examples --- is the price tag of
that asymmetry, and is reported in the same Table.
Both effects come from the same approximation : there is no free
lunch here, only a trade-off whose sign is exactly the one that
matches how prediction markets misprice tournaments.

\subsection{Prediction-market context}

The prediction-market angle is the practical motivation. Prediction
markets in general \cite{wolfers2004prediction} and their
combinatorial extensions \cite{hanson2003combinatorial} are well
understood as price-discovery instruments, but their bracket-level
implementations rarely encode the structural correlations between
matches. Polymarket crossed \$44~B in 2025 trading volume on its way
to an ICE-backed \$9~B valuation, and its futures pricing aggregates
over historical marginal frequencies rather than over
bracket-structural joint probabilities. A model that puts higher prior weight on
structurally-correlated upset paths would, ex ante, identify the
mispricings on the underdog side of that liquidity profile --- which
is also the side where the spreads are widest and the inventory is
shortest. That is the asymmetry the seeded walk exploits.

\subsection{A naive question, posed plainly}
\label{sec:candide}

One can pose the question naively : can a spectator watching a match
be quantum-mechanically entangled with the ball, the referee, or the
player at the moment of the goal ? The operational answer is no.
Macroscopic decoherence settles it : an object the size of a
football loses its quantum coherence in roughly $10^{-23}$ seconds
at room temperature \cite{joos_zeh_1985}, and the human brain,
energetically, is an even more hostile environment for coherence
\cite{tegmark_2000}. The question, taken literally, is closed.

One notes nonetheless, on re-reading the encoding of
Section~\ref{sec:method}, that the construction belongs to the same
mathematical family as objects studied for fifteen years in the
still-young field of quantum biology \cite{lambert2013quantum}. Our
quantum walk on the graph of matches of a tournament shares its
structure with the excitonic walk of the FMO photosynthetic complex
\cite{engel_2007, plenio2008dephasing}, where coherent energy packets
propagate on a graph of conditionally-coupled chromophores. The pair of entangled qubits
that encodes a conditional matchup in our construction shares its
structure with the pair of entangled radicals that migratory birds
sustain in their retinal cryptochrome for geomagnetic orientation
\cite{hore_mouritsen_2016}. The physical scales are separated by
thirteen orders of magnitude, but the mathematical object underneath,
a multipartite entangled state whose correlations are conditional
on a graph, is rigorously the same.

We push the parallel no further here. Our model remains a pricing
tool, not a physical substrate. That the same primitive appears
both at the femtosecond scale of energy transport in leaves and at
the hourly scale of a single-elimination bracket seems however
worth a mention, less as an argument for this paper than as a hint
at what this structure might price elsewhere.

\subsection{Extensions and follow-up work}

The Grover amplification of Section~\ref{sec:grover} and the QAE of
Section~\ref{sec:qae} together cover the simplest variant of the
path-likelihood family already flagged by the 2018 \emph{Quantum
Frontiers} note \cite{quantumfrontiers_wc_2018} : an unconditional
``team $X$ wins the tournament'' subspace, queried both by
amplification (which rotates $p_0$ to peak) and by amplitude
estimation (which returns $\hat p_0$ directly). The next natural
step is the conditional form ``given that $X$ reaches the semifinal,
who is the most likely champion ?'', which couples the same oracle
construction with an extra conditioning register and which we leave
to follow-up work. A faithful encoding of all $2^{r-1}$ parent
configurations would in parallel eliminate the seeded bias, at the
cost of circuit depth growing with bracket depth. NISQ-friendly QAE
variants without phase estimation \cite{suzuki2020amplitude} cut the
ancilla overhead substantially and are a natural target for near-term
hardware. All extensions remain tractable in simulation up to roughly
25 qubits on a single GPU via cuQuantum
\cite{cuquantum_benchmark_2023}.

\FloatBarrier
\section{Conclusion}

We have written down a quantum-walk encoding for bracket-structured
probability estimation, formalised the seeded approximation it
relies on at deeper rounds, and verified that the Qiskit Aer
implementation matches a closed-form enumeration of the same rule
to within sampling noise. At the current 15-qubit simulator scale
the contribution is conceptual rather than algorithmic : no quantum
advantage is claimed, but the encoding reproduces bracket
correlation structure that classical Monte Carlo only recovers
through explicit path conditioning.

The back-tests on five real 2025--2026 brackets are the main
empirical message. The walk upweights every realised underdog deep
run in our sample (NBA Spurs, RG WTA Andreeva and Chwalinska, RG ATP
Mensik and Arnaldi), with $Q/C$ ratios between $1.71\times$ and
$4.44\times$. It does not produce upset signal where the rating
spread does not support it (UCL PSG champion at $\approx 1.01$, NHL
Vegas--Colorado near-coin-flip at $0.63\times$, NBA OKC champion at
$1.04\times$). It does down-weight top seeds whose Elo lead is
moderate enough to be eroded by repeated conditional matchups
(NYK reach-Finals marginal at $0.71\times$, NYK champion at
$0.50\times$). All three behaviours come from the same approximation
and are reported with
the same Table. For prediction-market pricing the sign of the bias
is in the right direction : the under-pricing of structurally-
correlated upset paths is precisely the population the seeded
approximation preferentially identifies.

We additionally implement, on the same encoding, both Grover
amplitude amplification (Section~\ref{sec:grover}) and Quantum
Amplitude Estimation (Section~\ref{sec:qae}) for the SAS-champion
subspace of the NBA bracket. The amplification curve peaks at
$97.89\%$ on $k = 4$ iterations, and the QAE returns
$\hat p_{\text{mode}} = 3.81\%$ at $m = 6$ ancillas (21 qubits
total), within the $\pi/2^m$ resolution band of the reference
$p_0 = 3.46\%$ from unamplified sampling. Together the two queries
realise the canonical Brassard-Hoyer-Mosca-Tapp pipeline on a real
bracket, closing the most basic of the Grover-based extensions
originally flagged as future work by the 2018 \emph{Quantum
Frontiers} note.

Finally we issue, before the June~11 kickoff, a prospective
prediction for the FIFA World Cup 2026 R16-onward under a balanced
1-vs-16 bracket assumption. Mid-tier nations (NED, BEL, AUT, COL)
are upweighted by 2--3 percentage points over classical MC, and the
top favourites (FRA, ESP, BRA) reduced by a commensurate amount.
The encoding scales to roughly 25 qubits on a single GPU via
cuQuantum, and is a natural candidate for the late NISQ window and
the 50--100 logical-qubit regime expected from the first generations
of fault-tolerant hardware.

\bibliographystyle{plain}
\bibliography{references}

\end{document}
